% Präambel des Benutzerhandbuchs – gemeinsame Pakete und Makros.
% Übersetzen mit LuaLaTeX:  latexmk -lualatex handbuch.tex   (oder: make)

\usepackage{fontspec}
\usepackage[ngerman]{babel}
\usepackage[autostyle=true,german=quotes]{csquotes}
\usepackage{lmodern}
\usepackage{microtype}
\usepackage{amsmath,amssymb}
\usepackage{graphicx}
\usepackage[table]{xcolor}
\usepackage{booktabs}
\usepackage{longtable}
\usepackage{tabularx}
\usepackage{xltabular}
\usepackage{array}
\usepackage{enumitem}
\usepackage[locale=DE,per-mode=symbol]{siunitx}
\usepackage[os=win]{menukeys}
\usepackage[most]{tcolorbox}
\usepackage{listings}
\usepackage{imakeidx}
\usepackage[hidelinks,unicode,bookmarksnumbered]{hyperref}

\graphicspath{{../images/}{bilder/}}
\makeindex[intoc,columns=2,title=Stichwortverzeichnis]

\KOMAoptions{parskip=half}
\setlist{itemsep=2pt,topsep=4pt}
\setcounter{tocdepth}{1}

% ---------------------------------------------------------------- Farben
\definecolor{rtblau}{HTML}{3D5A80}
\definecolor{rtgelb}{HTML}{E0B040}
\definecolor{rtgrau}{HTML}{F2F3F5}
\definecolor{rtrot}{HTML}{B03A2E}
\definecolor{rtgruen}{HTML}{2E7D32}

% ---------------------------------------------------------------- Oberfläche
% \gui{Text}           Beschriftung eines Bedienelements (Knopf, Feld, Dock, Fenster)
% \menu{File,Save As...}  Menüpfad (Einträge durch Kommas getrennt)
% \keys{Strg+Z}        Tastenkombination (Teile durch + getrennt)
% \maus{rechts ziehen} Mausgeste
% \cmd{Line}           Befehl der Befehlszeile (wie getippt)
% \param{rays}         Name eines Plugin-Parameters / Schlüssel in Dateien
% \datei{a/b.rttf}     Datei- oder Ordnername
% \code{text}          kurzer Code / Eingabe
% \plugin{Raytracer}   Name eines Plugins
% \block{Camera}       Block des Node-Workflows
\newcommand{\gui}[1]{\textsf{\bfseries #1}}
\newcommand{\maus}[1]{\textit{#1}}
\providecommand{\cmd}{}% menukeys defines \cmd (Mac command key); not used here
\renewcommand{\cmd}[1]{\texttt{\textbf{#1}}\index{#1@\texttt{#1} (Befehl)}}
\newcommand{\param}[1]{\texttt{#1}}
\newcommand{\datei}[1]{\texttt{#1}}
\newcommand{\code}[1]{\texttt{#1}}
\newcommand{\plugin}[1]{\textsf{#1}}
\newcommand{\block}[1]{\textsf{\itshape #1}}
\newcommand{\produkt}{Raytatouille Fusion}
\newcommand{\version}{1.0}

% ---------------------------------------------------------------- Kästen
\newtcolorbox{hinweis}{enhanced,breakable,colback=rtgrau,colframe=rtblau,boxrule=0.5pt,
  left=6pt,right=6pt,top=4pt,bottom=4pt,title={Hinweis},fonttitle=\sffamily\bfseries\small,
  coltitle=white}
\newtcolorbox{tipp}{enhanced,breakable,colback=rtgruen!5,colframe=rtgruen,boxrule=0.5pt,
  left=6pt,right=6pt,top=4pt,bottom=4pt,title={Tipp},fonttitle=\sffamily\bfseries\small}
\newtcolorbox{warnung}{enhanced,breakable,colback=rtrot!5,colframe=rtrot,boxrule=0.5pt,
  left=6pt,right=6pt,top=4pt,bottom=4pt,title={Achtung},fonttitle=\sffamily\bfseries\small}

% Schritt-für-Schritt-Anleitungen
\newlist{schritte}{enumerate}{1}
\setlist[schritte]{label=\arabic*.,leftmargin=*,itemsep=3pt}

% Abbildung: \abbildung[Breite]{datei}{Unterschrift}{label}
\newcommand{\abbildung}[4][\linewidth]{%
  \begin{figure}[htbp]\centering\includegraphics[width=#1]{#2}\caption{#3}\label{#4}\end{figure}}

% Listings für Dateien und Befehle im Terminal
\lstdefinestyle{terminal}{basicstyle=\ttfamily\small,backgroundcolor=\color{rtgrau},
  frame=single,framerule=0pt,breaklines=true,columns=fullflexible,keepspaces=true,
  xleftmargin=4pt,xrightmargin=4pt,aboveskip=6pt,belowskip=6pt}
\lstset{style=terminal}

% Referenztabellen: \begin{referenz}{Spaltenformat} ... \end{referenz}  (seitenübergreifend)
\newenvironment{referenz}[1]{\small\xltabular{\linewidth}{#1}}{\endxltabular}
